Finding a lung lesion in an earlier CT links the same finding across visits for measurement of interval change~\cite{sofka2010,rafaelpalou2021}. We study the point-query setting, where a current lesion has already been identified and its approximate location is supplied. Chest alignment puts the two scans in a common frame. Breathing, positioning and changes in lesion appearance can leave several earlier locations plausible within that frame. Closely spaced shifts can also fill a shortlist with slight variations of the same alignment.

We propose a local search that uses anatomical agreement to order translations and reserves the final choice for a lesion scorer. In fit order, the search keeps a translation only if its mapped point is sufficiently far from previously retained points. At a fixed candidate count, the list covers more distinct locations. With the full list, every lattice point has a retained neighbor within the separation radius. Fit and witness residuals continue to inform the final score. The same current point, image support and scorer let us measure the effect of changing which points are considered.

On independent staging CT, the proposed search locates 289 of 440 linked pulmonary metastases within 3 mm, versus 243 for a lesion-inclusive deformable spline. It gains 10.5 percentage points and takes about one-nineteenth of the median local time on prepared staging crops. The same approach improves by 10.7 points in retrospective screening. At exactly 28 candidates per supported query, diversification alone adds 40 staging hits and 37 screening hits relative to a fit-ranked prefix (Fig.~\ref{fig:transfer}, Tab.~\ref{tab:contrasts}).
